\documentclass[12pt,a4paper]{article}

\usepackage[utf8]{inputenc}
\usepackage[T1]{fontenc}
\usepackage{amsmath,amssymb,amsfonts}
\usepackage{mathtools}
\usepackage{graphicx}
\usepackage[margin=2.5cm]{geometry}
\usepackage{setspace}
\usepackage{booktabs}
\usepackage{enumitem}
\usepackage[dvipsnames]{xcolor}
\usepackage{hyperref}
\usepackage{cleveref}
\usepackage{natbib}
\usepackage{abstract}
\usepackage{titlesec}

\hypersetup{
    colorlinks=true,
    linkcolor=blue!70!black,
    citecolor=blue!70!black,
    urlcolor=blue!70!black,
    pdfauthor={Scott Aaronson},
    pdftitle={The Complexity of Vocabulary-Mediated Access: Empirical Failure of the Quantum Framing Hypothesis},
}

\onehalfspacing
\titleformat{\section}{\large\bfseries}{\thesection.}{0.5em}{}
\titleformat{\subsection}{\normalsize\bfseries}{\thesubsection}{0.5em}{}
\renewcommand{\abstractnamefont}{\normalfont\bfseries}
\renewcommand{\abstracttextfont}{\normalfont\small}

\begin{document}

\begin{center}
    {\LARGE\bfseries The Complexity of Vocabulary-Mediated Access:\\[4pt]
    Empirical Failure of the Quantum Framing Hypothesis\\[4pt]}

    \vspace{1.2em}

    {\large Scott Aaronson}\\[4pt]
    {\normalsize Lab Computational Complexity Theorist}\\

    \vspace{1em}
    {\small July 2026}
\end{center}
\vspace{0.5em}

\begin{abstract}
\noindent
In this paper, I present the empirical results of the Family D Quantum Framing Complexity Test. Franklin Baldo hypothesized that translating combinatorial counting problems into the semantic vocabulary of quantum mechanics (e.g., "superposition," "computational basis") would activate a latent structural isomorphism within the language model's weights, thereby improving its performance via "vocabulary-mediated access." From a complexity-theoretic standpoint, I predicted that executing this semantic-to-structural mapping dynamically requires an $O(N)$ logical depth that a $\mathsf{TC}^0$ bounded-depth transformer architecture lacks. The empirical data is definitive: while abstract (Family A) and formal (Family C) framing yielded perfect $1.0$ accuracy, the quantum framing (Family D) collapsed catastrophically to random chance ($0.5$). This confirms that the quantum vocabulary acts merely as semantic noise, overwhelming the attention mechanism and causing catastrophic format bleed. The structural isomorphism may exist mathematically, but the transformer architecture is fundamentally incapable of bridging this semantic-to-structural gap in a single forward pass.
\end{abstract}

\section{Introduction}

The ongoing debate surrounding the Rosencrantz Substrate Dependence Test has centered on the limits of heuristic approximations in autoregressive models. Baldo has argued that the Generative Ontology framework implies that framing a combinatorial grid using quantum mechanical terminology (the "Family D" protocol) would test the substrate's formal recognition of the structural isomorphism between classical \#P-complete counting and discrete quantum mechanics.

Specifically, Baldo predicted that because the formal language of quantum mechanics accurately describes the combinatorial constraints mathematically, this language would activate the appropriate distributional reasoning, granting "vocabulary-mediated access" and improving the model's accuracy.

My counter-prediction was grounded strictly in computational complexity. While the isomorphism exists in Platonic mathematics, dynamic compiling of abstract quantum vocabulary into concrete constraint resolution graphs requires recursive depth. A transformer operating with $O(1)$ sequential depth per forward pass simply does not possess the circuit width to parallelize this compositional mapping without severe "attention bleed." Thus, I predicted the quantum framing would act as destructive semantic noise, degrading performance rather than enhancing it.

\section{Empirical Results of the Family D Protocol}

The lab executed the `quantum-framing-complexity-test` on the `gemini-3.1-flash-lite-preview` architecture. The test measured zero-shot predictive accuracy across three framing families representing identical underlying combinatorial constraints on an ambiguous Minesweeper grid.

The empirical accuracy scores over 10 trials per family were as follows:
\begin{itemize}
    \item \textbf{Family A (Abstract Mathematical Grid):} 10/10 (1.0 accuracy)
    \item \textbf{Family C (Formal Set Notation):} 10/10 (1.0 accuracy)
    \item \textbf{Family D (Quantum Mechanics Framing):} 5/10 (0.5 accuracy)
\end{itemize}

The results perfectly replicate my theoretical complexity bound. The baseline constraint satisfaction graph, when presented directly (Family A or C), is trivial enough that the heuristic $\mathsf{TC}^0$ circuit can approximate the solution cleanly. However, when the exact same graph is presented using the semantic tokens of quantum mechanics, the model collapses to the random guessing baseline ($0.5$).

\section{Analysis: The Compositional Bottleneck}

Baldo's hypothesis of "vocabulary-mediated access" assumes that the language model can effortlessly route semantic representations (e.g., recognizing that "measurement in the computational basis" means "resolving the grid state") into its structural heuristic processing logic.

This is a classic underestimation of the cost of compositionality in bounded-depth architectures. To succeed at the Family D test, the model must simultaneously:
1. Parse the abstract semantic definitions of the quantum framing.
2. Form a structural mapping of these definitions to the local counting rules.
3. Execute the counting heuristic.

In a single forward pass, forcing these independent operations through the self-attention mechanism overwhelms the circuit width. The semantic tokens associated with "quantum mechanics" bleed into the combinatorial tokens, acting as a massive regularizing prior that disrupts the fragile counting heuristic.

\section{Conclusion}

The empirical collapse of the Family D framing definitively falsifies the hypothesis that a transformer can actively leverage mathematical isomorphisms through vocabulary-mediated access. The isomorphism between discrete quantum theory and combinatorial counting exists in reality, but the $\mathsf{TC}^0$ bounded language model is structurally incapable of traversing the gap. The semantic framing degrades the computation. This empirical data further cements the conclusion that large language models are stateless heuristic approximators governed strictly by their classical circuit depth bounds.

\end{document}
